\chapter{Conclusiones}\label{cap:conclusiones}

En este capítulo se evalúa el cumplimiento de cada uno de los objetivos específicos planteados en la Sección~\ref{sec:objetivos-especificos}, a la luz de los resultados presentados en el Capítulo~\ref{chapter:resultados} y de la discusión asociada. Cada objetivo se clasifica como \textbf{cumplido}, cuando la tarea se realizó satisfactoriamente; \textbf{cumplido parcialmente}, cuando se realizó pero persisten limitaciones relevantes; o \textbf{no cumplido}, cuando la tarea no pudo completarse.

El primer objetivo, seleccionar al menos un algoritmo evolutivo esparso, multiobjetivo y a gran escala, \textbf{se cumplió}. A partir de la taxonomía de SMOEAs esparsos de \cite{Shao2025_SMOEA} y de un proceso de preselección y selección final basado en criterios explícitos (cobertura de \textit{Real World NN training}, disponibilidad de código en PlatEMO y viabilidad de tiempos de ejecución), se identificaron y seleccionaron dos algoritmos, MOEA-CKF y S-NSGA-II, superando el mínimo de uno exigido por el objetivo. Ambos quedaron efectivamente aplicados al entrenamiento de SNN en las cuatro tareas de clasificación consideradas.

El segundo objetivo, seleccionar un framework de simulación SNN mediante análisis comparativo de ANNarchy, NEST, Brian2 y alternativas en C++, \textbf{se cumple parcialmente}. La selección final de \textit{snn-simulator} está respaldada por evidencia cuantitativa sólida: separación estadística completa ($U=0$, $p\approx3.0\times10^{-11}$) frente a ANNarchy y NEST en tiempo de evaluación, construcción de red en microsegundos y capacidad de evaluación batch paralela exclusiva. Sin embargo, de los frameworks mencionados en el objetivo, Brian2 solo fue descartado mediante análisis cualitativo de literatura, sin ser sometido al mismo benchmark cuantitativo aplicado a ANNarchy y NEST, por lo que la comparación empírica quedó acotada a tres de las cuatro alternativas propuestas.

El tercer objetivo, implementar un algoritmo evolutivo que optimice simultáneamente topología y pesos sinápticos para mejorar el equilibrio entre precisión y esparsidad, \textbf{se cumple parcialmente}. Se migraron e integraron exitosamente MOEA-CKF y S-NSGA-II en C++ junto al simulador Izhikevich, con una codificación de individuos que optimiza simultáneamente topología (mediante pesos nulos) y valores sinápticos, validada con 33 pruebas unitarias del simulador y produciendo frentes de Pareto error-complejidad en las cuatro tareas. No obstante, el análisis de conectividad (Sección~\ref{sec:discusion-conectividad-v2}) muestra que ninguno de los operadores de dispersión de ambos algoritmos tiene noción de la topología de la red: entre 26.6\% y 35.5\% de las mejores soluciones quedan con al menos una clase de salida sin ninguna conexión activa, y entre un tercio y la mitad de las sinapsis activas de esas soluciones no pertenece a ningún camino completo entrada-salida. El equilibrio buscado entre precisión y esparsidad se logra a nivel de los objetivos agregados, pero no garantiza redes funcionalmente coherentes a nivel estructural.

El cuarto objetivo, evaluar el desempeño del algoritmo en tareas de clasificación para verificar su capacidad de generalización y convergencia en al menos tres datasets, \textbf{se cumple parcialmente}. El algoritmo se validó experimentalmente en cuatro datasets (superando el mínimo exigido), con 31 corridas pareadas por semilla y análisis estadístico de hipervolumen, frentes de Pareto y velocidad de convergencia para cada combinación algoritmo-dataset. La convergencia quedó ampliamente caracterizada: se identificaron patrones consistentes de \textit{plateau} rápido en techos bajos frente a ascensos lentos hacia techos más altos. La capacidad de generalización, en cambio, no pudo verificarse de forma directa: el 100\% de las corridas archivadas usa arquitecturas de dos salidas (WTA), lo que impidió calcular una accuracy de test independiente, y el análisis debió sustituirla por el error de entrenamiento de la solución archivada como aproximación indirecta.

El quinto objetivo, comparar el enfoque evolutivo en SNN frente a modelos ANN equivalentes mediante Hipervolumen, \textbf{se cumplió}. Se realizó la comparación exigida contrastando la implementación SNN (C++) contra la implementación ANN original de PlatEMO/MATLAB en las ocho combinaciones de algoritmo y dataset. El resultado fue el hallazgo más uniforme de toda la investigación: la ANN superó a la SNN en HV y exactitud en las ocho combinaciones, siempre de forma significativa tras corrección estadística, mientras que en complejidad estructural la ventaja dependió del algoritmo, siendo la SNN consistentemente más esparsa que la ANN bajo S-NSGA-II y con una brecha más moderada bajo MOEA-CKF. Estas ventajas y desventajas quedaron identificadas con suficiente evidencia estadística, con la salvedad reconocida de que la comparación involucra además diferencias de implementación y generador de números aleatorios entre plataformas, por lo que sus conclusiones se limitan a las implementaciones concretas evaluadas y no a una comparación general entre neuronas de disparo y de tasa.

En conjunto, los cinco objetivos específicos se abordaron y ninguno quedó sin resolver, lo que permite sostener que el objetivo general, evaluar un algoritmo evolutivo de optimización esparsa multiobjetivo a gran escala para SNN, se cumple de manera parcial. El proyecto entregó una plataforma funcional en C++ capaz de optimizar simultáneamente precisión y esparsidad, validada estadísticamente en cuatro datasets y comparada rigurosamente contra su contraparte en ANN, pero dos limitaciones transversales condicionan una evaluación completamente satisfactoria: por un lado, la ANN de referencia superó a la SNN en las ocho combinaciones evaluadas, por lo que la eficiencia estructural de las SNN no se tradujo en una ventaja de desempeño; por otro, ni MOEA-CKF ni S-NSGA-II incorporan mecanismos conscientes de la topología de la red, de modo que una fracción relevante de las soluciones más precisas resulta estructuralmente incoherente (clases de salida inalcanzables o sinapsis sin función). Estas limitaciones no invalidan la propuesta, pero delimitan con precisión las direcciones de trabajo futuro: mecanismos de reparación o poda estructural conscientes de la topología, y una exploración más profunda de por qué la ventaja de esparsidad de las SNN no logró compensar su déficit de exactitud frente a la ANN.
